\section{Discussion and Limitations}
\label{sec:discussion}
\paragraph{Availability and complete recovery.}
The Raw/Raw+E comparison shows that supplying explicit relations is associated with much stronger recovery of the declared entity-route reference on the matched archived tasks. Native metadata may omit the required links, so low Raw recall cannot by itself establish a reasoning error on sufficient evidence. The remaining errors under explicit relations pose a different question: whether the model retains the complete set specified by the query. High precision with incomplete recall shows that a mostly reference-consistent selection can still omit required route entities.

\paragraph{Interpreting structural profiles.}
Task-family differences remain visible at common sizes, while size profiles vary by family and model. These observations motivate reporting selection performance by query family and evidence condition. They do not establish a universal motif ranking or an isolated topology coefficient: families also differ in traversal direction, stopping rules, relation types, and sampled entities. Different sizes use different cohorts, and the fixed catalogue changes target prevalence and distractor count as reference size increases.

\paragraph{Presentation and residual errors.}
Repeated first-neighborhood outputs provide a localized description of some incoming-tree errors. Their paired transitions show why a marginal reduction in an error signature is insufficient: an output may become exact, remain incomplete, or change to another error. These are observations of final sets, not evidence of hidden search depth or internal stopping. Their detailed scope is the inspected Core incoming-tree slice rather than all families and sizes.

Both-card hosting motivates a hypothesis about access to distributed facts, but also changes repetition, salience, and input length. Destination-only hosting does not produce the same pattern across models. Separating these explanations would require controls that hold repetition and length fixed while varying hosting. Likewise, adjacent-context differences can involve distractor identities and graph density as well as adjacency. The archived comparisons do not identify a unique attention, traversal, or discrimination mechanism.

\paragraph{What the selection objective identifies.}
We score the union of allowed route vertices, including the anchor. This is distinct from minimal graph support, indispensable visible cards, and correct endpoint answers. A smaller selection may suffice for an endpoint answer even when it fails the route-coverage convention. High recall alone also permits selecting everything, which has recall one by construction; precision and exact are therefore necessary companion diagnostics.

Some rule violations are observationally indistinguishable in simple contexts. On the 100 inspected Qwen K1 six-node incoming-tree inputs, ignoring direction or increasing the depth bound from two to three produces the same coverage set as the correct rule. Correct final sets in that slice cannot establish adherence to those rules. Adjacent contexts can supply distinguishing cases, but different rules have different diagnostic support. A full anchor-connected-component baseline has not been measured; neither ideal performance by that heuristic nor the absence of shortcuts is claimed.

\paragraph{Scope and reproducibility.}
The study covers one reused corporate KG~\cite{kirpichev2026towards} and frozen model/protocol labels, with no generation-seed replication or new KG evaluation. It reanalyzes existing tasks and saved outcomes; it does not introduce a newly collected dataset or independent experimental replication. Base tasks, anchors, and reference entities overlap, so descriptive averages do not remove dependence. We introduce no recall/precision significance or population-causal claims.

Frozen-reference checks establish consistency of the presented tasks, not completeness against the unloaded source database. Missing Gemma replies/settings, native sufficiency, and incremental histories remain limitations. Saved rows, selectors, denominators, hashes, and plotting scripts support reproduction of these summaries; resource redistribution and graph access are separate matters.

\paragraph{Responsibility and downstream scope.}
This study does not evaluate a new human-subject population or establish downstream QA improvement. Generative AI assistance supported evidence organization, plotting code, and manuscript drafting; experimental values derive from saved outcomes. Authors remain responsible for source verification, anonymity, resource permissions, and the provenance of reused materials.
